\begin{ex}
Tính: $2 \times 8 = \dots$
\shortans{16}
\loigiai{
$2 \times 8 = 16$.
}
\end{ex}

\begin{ex}
Tính: $5 \times 7 = \dots$
\shortans{35}
\loigiai{
$5 \times 7 = 35$.
}
\end{ex}

\begin{ex}
Tính: $18 : 2 = \dots$
\shortans{9}
\loigiai{
$18 : 2 = 9$.
}
\end{ex}

\begin{ex}
Tính: $45 : 5 = \dots$
\shortans{9}
\loigiai{
$45 : 5 = 9$.
}
\end{ex}

\begin{ex}
Tính: $2 \times 9 + 15 = \dots$
\shortans{33}
\loigiai{
$2 \times 9 + 15 = 18 + 15 = 33$.
}
\end{ex}

\begin{ex}
Tính: $5 \times 8 - 18 = \dots$
\shortans{22}
\loigiai{
$5 \times 8 - 18 = 40 - 18 = 22$.
}
\end{ex}

\begin{ex}
Tính: $20 : 2 + 35 = \dots$
\shortans{45}
\loigiai{
$10 + 35 = 45$.
}
\end{ex}

\begin{ex}
Tính: $50 : 5 - 6 = \dots$
\shortans{4}
\loigiai{
$10 - 6 = 4$.
}
\end{ex}

\begin{ex}
Mỗi đôi đũa có 2 chiếc đũa. Hỏi 9 đôi đũa có bao nhiêu chiếc đũa?
\shortans{18}
\loigiai{
$2 \times 9 = 18$ chiếc.
}
\end{ex}

\begin{ex}
Có 16 học sinh chia đều thành 2 nhóm. Hỏi mỗi nhóm có bao nhiêu học sinh?
\shortans{8}
\loigiai{
$16 : 2 = 8$ học sinh.
}
\end{ex}

\begin{ex}
Mỗi tuần lễ mẹ đi làm 5 ngày. Hỏi trong 6 tuần lễ mẹ đi làm tất cả bao nhiêu ngày?
\shortans{30}
\loigiai{
$5 \times 6 = 30$ ngày.
}
\end{ex}

\begin{ex}
Có 35 lít dầu chia đều vào các can, mỗi can chứa 5 lít. Hỏi rót được bao nhiêu can dầu?
\shortans{7}
\loigiai{
$35 : 5 = 7$ can.
}
\end{ex}

\begin{ex}
Có 20 cái kẹo chia đều cho 5 bạn. Hỏi mỗi bạn được mấy cái kẹo?
\shortans{4}
\loigiai{
$20 : 5 = 4$ cái kẹo.
}
\end{ex}

\begin{ex}
Điền số thích hợp: $2 \times \dots = 14$.
\shortans{7}
\loigiai{
$14 : 2 = 7$.
}
\end{ex}

\begin{ex}
Điền số thích hợp: $5 \times \dots = 40$.
\shortans{8}
\loigiai{
$40 : 5 = 8$.
}
\end{ex}

\begin{ex}
Điền số thích hợp: $\dots : 2 = 6$.
\shortans{12}
\loigiai{
$6 \times 2 = 12$.
}
\end{ex}

\begin{ex}
Điền số thích hợp: $\dots : 5 = 5$.
\shortans{25}
\loigiai{
$5 \times 5 = 25$.
}
\end{ex}

\begin{ex}
Lớp 3A có 5 tổ, mỗi tổ có 6 bạn. Hỏi lớp 3A có bao nhiêu bạn?
\shortans{30}
\loigiai{
$5 \times 6 = 30$ bạn.
}
\end{ex}

\begin{ex}
Mỗi bàn học có 2 học sinh ngồi. Có tất cả 10 bàn học. Hỏi có bao nhiêu học sinh ngồi?
\shortans{20}
\loigiai{
$2 \times 10 = 20$ học sinh.
}
\end{ex}

\begin{ex}
Tính: $2 \times 6 : 2 = \dots$
\shortans{6}
\loigiai{
$12 : 2 = 6$.
}
\end{ex}

\begin{ex}
Tính: $5 \times 4 : 2 = \dots$
\shortans{10}
\loigiai{
$20 : 2 = 10$.
}
\end{ex}

\begin{ex}
Một con gà có 2 chân. Hỏi 8 con gà có tất cả bao nhiêu chân?
\shortans{16}
\loigiai{
$2 \times 8 = 16$ chân.
}
\end{ex}

\begin{ex}
Mỗi bàn tay có 5 ngón tay. Hỏi 2 bàn tay có bao nhiêu ngón tay?
\shortans{10}
\loigiai{
$5 \times 2 = 10$ ngón tay.
}
\end{ex}

\begin{ex}
Có 14 quả cam xếp đều vào 2 đĩa. Hỏi mỗi đĩa có bao nhiêu quả cam?
\shortans{7}
\loigiai{
$14 : 2 = 7$ quả cam.
}
\end{ex}

\begin{ex}
Có 30 bông hoa cắm đều vào 5 lọ. Hỏi mỗi lọ có bao nhiêu bông hoa?
\shortans{6}
\loigiai{
$30 : 5 = 6$ bông hoa.
}
\end{ex}
